% 02-problem.tex — §2 Background and motivation
\section{Background and Motivation}
\label{sec:problem}

\subsection{Hardware Cost and Carry-Chain Shape of U256 Arithmetic}

The native word length of the EVM is 256 bits, whereas general-purpose
registers on current commodity hardware are 64 bits wide. An EVM-level
U256 addition expands, after JIT code generation, into four machine-level
additions: one ordinary 64-bit \texttt{add} followed by three
carry-propagating \texttt{adc} instructions that form a \emph{serial
carry chain}. Subtraction is analogous, using one \texttt{sub} and three
\texttt{sbb} instructions. The listing below shows the 4-limb code on
both sides for a typical U256 addition:

\begin{center}
\small
\begin{tabular}{@{}l|l@{}}
\textbf{dMIR (4-limb addition)} & \textbf{x86-64 output} \\
\hline
\texttt{r0 = add(a0, b0)}   & \texttt{add rax, rcx}  \\
\texttt{r1 = ADC(a1, b1, cout0)} & \texttt{adc rbx, rdx}  \\
\texttt{r2 = ADC(a2, b2, cout1)} & \texttt{adc rsi, rdi}  \\
\texttt{r3 = ADC(a3, b3, cout2)} & \texttt{adc r8,\enspace r9} \\
\end{tabular}
\end{center}

Among all instructions emitted by the JIT, instructions directly tied to
U256 multi-word arithmetic account for about \textbf{2.15\%} (weighted
mean across the U256-heavy subset of the 27 benchmarks in
\S\ref{sec:evaluation}---the leftmost cluster ``U256-heavy (1--11)'' in
Figure~\ref{fig:speedup}). The share
is small, but the carry chain has a special property: every
\texttt{adc}/\texttt{sbb} on the chain depends on the carry output of the
previous instruction, forming a strict \emph{serial data dependency}
that instruction-level parallelism cannot hide.
U256 instructions are therefore few in number yet contribute
disproportionately to execution time.
The family-attribution data in \S\ref{sec:evaluation} show that on
U256-heavy benchmarks the U256 rule hit rate ranges from 25\% to 99.9\%
with a very uneven distribution; on non-U256-heavy benchmarks (sha1
family, \texttt{jump\_around}) it is well below that range and gains
come from x86-layer mechanical cleanup.

\subsection{Expressiveness Limit of General-Purpose IR on the Carry Chain}
\label{sec:expressiveness}

In a general-purpose compiler IR such as LLVM, the basic arithmetic
operations (\texttt{add}, \texttt{sub}, \texttt{mul}) are defined only on
$N$-bit integers \texttt{iN}. To obtain the carry bit explicitly, the
programmer must call the intrinsic
\texttt{@llvm.uadd.with.overflow.iN}, which returns a \texttt{\{iN, i1\}}
multi-return tuple---the result and the carry bit are packed into the
same structure. The carry bit is therefore \emph{not a native type} in
LLVM IR but is hidden in the second field of a struct.

This design limits the expressiveness of automatic equivalence checkers.
Following the peephole verification framework proposed by Lopes et al.\
in Alive~\cite{alive} and inherited by Alive2~\cite{alive2},
Table~\ref{tab:cannot-express} lists four typical carry-chain rules that
Alive2 cannot directly verify for equivalence---Alive2 returns a
type mismatch on the multi-return tuple or times out on the lifted
encoding. Under Alive2 these
four rules fall in the \textbf{not automatically verifiable} category
(within the supported encoding); the limitation stems from how the
general-purpose IR models the carry bit and is independent of the rule
implementation and the solver capacity.

% tables/tab2-1-cannot-express.tex — 4 carry-chain rules that general-purpose IR cannot verify directly
\begin{table}[t]
\centering
\caption{Four typical carry-chain rules that Alive2 (based on LLVM IR)
cannot directly verify for equivalence because of the multi-return
tuple encoding. The first three are single \texttt{ADC}/\texttt{SBB}
rules where the right-hand side drops the carry output; the fourth is a
4-limb carry-chain merge rule.}
\label{tab:cannot-express}
\footnotesize
\begin{tabular}{@{}ll@{}}
\toprule
Rule & Rewrite \\
\midrule
\texttt{adc\_zero\_carry}  & $\texttt{ADC}(x,y,0) \to \texttt{add}(x,y)$ \\
\texttt{sbb\_zero\_borrow} & $\texttt{SBB}(x,y,0) \to \texttt{sub}(x,y)$ \\
\texttt{sbb\_self\_zero}   & $\texttt{SBB}(x,x,0) \to 0$ \\
\midrule
4-limb ADC chain & $4 \times \texttt{uadd.w.overflow} \to \texttt{add.i256}$ \\
\bottomrule
\end{tabular}
\end{table}


The Alive2 command lines, failure messages, and version numbers
(Alive2 mainline v21.0, Z3 4.8.12) for all four rules are archived
in the artifact
(\S\ref{sec:availability}); the relation to Hydra~\cite{hydra}
(OOPSLA 2024, a different system from the Hydra Framework at USENIX
Security 2018 cited in \S\ref{sec:related}) is discussed in
\S\ref{sec:related}.

The bottleneck is not the capacity of the SMT
solver~\cite{demoura2008z3}. The remaining six U256 rules of the same
class, which do not involve multi-return carry propagation, all pass
Alive2 equivalence checking with an average runtime of about
60\,ms. The 256-bit bit-vector size is not an obstacle for the solver;
the bottleneck lies in how the IR expresses the carry bit.

Concretely, expressing carry-propagating addition in LLVM IR requires
five lines around \texttt{uadd.with.overflow} to extract the result and
carry bit from a multi-return tuple, whereas a single \texttt{ADC}
suffices in dMIR. Native carry-bit support brings the four rules above
into the scope of automatic equivalence checking; dMIR trades LLVM
interoperability for automatic SMT verification of carry-chain rules,
and Z3 discharges all four rules of Table~\ref{tab:cannot-express}
directly under the dMIR bit-vector model.
